% Part: applied-modal-logics
% Chapter: epistemic-logic
% Section: bisimulations

\documentclass[../../../include/open-logic-section]{subfiles}

\begin{document}

\olfileid[id]{aml}{el}{bsd}

\olsection{Bisimulasi}

Satu pertanyaan yang masih dapat kita ajukan mengenai daya ungkap bahasa
epistemik kita menyangkut hubungan antara model dan formula yang berlaku di
dalamnya. Dari hasil korespondensi kerangka, kita telah melihat bahwa jika
formula tertentu valid dalam suatu kerangka, formula itu juga memastikan
bahwa kerangka tersebut memenuhi sifat tertentu. Namun, apakah bahasa modal
kita, misalnya, memungkinkan kita membedakan suatu dunia yang mempunyai
panah refleksif dari rantai dunia tak hingga yang masing-masing menuju dunia
berikutnya? Dengan kata lain, adakah formula~$A$ yang berlaku hanya pada salah
satu dari kedua dunia tersebut?

Bisimulasi merupakan relasi yang dapat kita definisikan di antara model
relasional untuk menyatakan bahwa model-model itu pada dasarnya mempunyai
struktur yang sama. Seperti akan kita lihat, bisimulasi menangkap suatu
pengertian kesetaraan antarmodel yang dapat dinyatakan dalam bahasa epistemik
kita.

\begin{defn}[Bisimulasi]
Misalkan $M_1=\tuple{W_1,R_1,V_1}$ dan $M_2=\tuple{W_2,R_2,V_2}$ merupakan
dua model relasional. Misalkan pula $\mathcal{R}\subseteq W_1\times W_2$
merupakan relasi biner. Kita mengatakan bahwa $\mathcal{R}$ merupakan
\emph{bisimulasi} jika, untuk setiap
$\tuple{w_1,w_2}\in\mathcal{R}$, berlaku:

\begin{enumerate}
  \item $w_1\in V_1(p)$ jika dan hanya jika $w_2\in V_2(p)$ untuk semua
    !!{propositional variable}s~$p$.
  \item Untuk semua agen $a\in G$ dan dunia $v_1\in W_1$, jika
    $R_{1,a}w_1v_1$, maka terdapat $v_2\in W_2$ sedemikian sehingga
    $R_{2,a}w_2v_2$ dan $\tuple{v_1,v_2}\in\mathcal{R}$.
  \item Untuk semua agen $a\in G$ dan dunia $v_2\in W_2$, jika
    $R_{2,a}w_2v_2$, maka terdapat $v_1\in W_1$ sedemikian sehingga
    $R_{1,a}w_1v_1$ dan $\tuple{v_1,v_2}\in\mathcal{R}$.
\end{enumerate}

Jika terdapat bisimulasi antara $M_1$ dan $M_2$ yang menghubungkan dunia
$w_1$ dan $w_2$, kita juga dapat menulis
$\tuple{M_1,w_1}\leftrightarroweq\tuple{M_2,w_2}$ dan menyebut
$\tuple{M_1,w_1}$ dan $\tuple{M_2,w_2}$ \emph{bisimilar}.
\end{defn}

Klausa yang berbeda dalam relasi bisimulasi menjamin hal yang berbeda.
Klausa~1 menjamin bahwa dunia bisimilar memenuhi formula bebas modal yang
sama karena klausa itu menjamin kesesuaian pada semua !!{propositional
variable}s. Kedua klausa lainnya, yang masing-masing kadang disebut klausa
``maju'' dan ``mundur,'' menjamin bahwa relasi aksesibilitas mempunyai
struktur yang sama.

\begin{thm}
Jika $\tuple{M_1,w_1}\leftrightarroweq\tuple{M_2,w_2}$, maka, untuk setiap
!!{formula}~$!A$, berlaku $\mSat{M_1}{!A}[w_1]$ jika dan hanya jika
$\mSat{M_2}{!A}[w_2]$.
\end{thm}

\begin{figure}
  \begin{center}
    \begin{tikzpicture}[modal]
      \node[world] (w1) {$w_1$};
      \node[world] (w2) [above left=of w1]{$w_2$};
      \node[world] (w3) [above right=of w1] {$w_3$};
      \draw[<->] (w1) to node {$a$} (w2);
      \draw[->,loop below] (w1) to node {$a$} (w1);
      \draw[<->] (w1) to [swap] node {$a$} (w3);
      \draw[->,loop above] (w2) to node {$a$} (w2) ;
      \draw[->,loop above] (w3) to node {$a$} (w3) ;

      \node[world](v1)[right of=w1, xshift=1.5in]{$v_1$};
      \node[world](v2)[above of=v1]{$v_2$};
      \draw[<->] (v1) to node {$a$} (v2);
      \draw[->,loop below] (v1) to node {$a$} (v1);
      \draw[->,loop above] (v2) to node {$a$} (v2) ;

      \draw[-,dotted] (w1) to (v1) ;
      \draw[-,dotted,bend left=45] (w2) to (v2) ;
      \draw[-,dotted,bend left=30] (w3) to (v2) ;
    \end{tikzpicture}
  \end{center}
  \caption{Dua model bisimilar.}
  \ollabel{fig:bisimilar}
\end{figure}

Walaupun kedua model dalam \olref{fig:bisimilar} tidak sama persis, terdapat
bisimulasi yang menghubungkan dunia $w_1$ dan~$v_1$. Bisimulasi ini juga
menghubungkan $w_2$ dan $w_3$ dengan~$v_2$; gagasannya ialah bahwa tidak ada
sesuatu yang dapat dinyatakan dalam bahasa modal kita yang benar-benar dapat
membedakan dunia-dunia itu. Namun, situasinya akan berbeda jika $w_2$ dan
$w_3$ memenuhi !!{propositional variable}s yang berbeda.

\end{document}
